% Scope of the Work
\section{Scope of the Work}

This monograph addresses a specific and circumscribed question: \emph{What mathematical and logical operations can be demonstrated in the Indus script based on distributional, combinatorial, and structural evidence?}

\subsection{What This Work IS}

This is a formal reconstruction of the operational capabilities of the Indus script system, demonstrating:

\begin{enumerate}
    \item The script's capacity to encode finite arithmetic operations
    \item Its implementation of relational and logical structures
    \item Its function as a state-based validation system
    \item Its formal equivalence to deterministic finite automata
\end{enumerate}

The reconstruction is grounded in:
\begin{itemize}
    \item Empirical sign frequency data
    \item Distributional constraints and positional patterns
    \item Combinatorial analysis of valid sequences
    \item Cross-cultural comparison with functionally similar systems
\end{itemize}

\subsection{What This Work IS NOT}

This work explicitly does NOT:

\begin{enumerate}
    \item Propose a phonetic decipherment of the Indus script
    \item Claim to ``translate'' Indus inscriptions into natural language
    \item Identify the language(s) potentially represented by the script
    \item Assert that the script is or is not a ``true writing system'' (a distinction we find unhelpful)
    \item Provide semantic interpretations of individual signs beyond their operational function
\end{enumerate}

\subsection{Methodological Position}

We adopt a functional-operational approach: the script is analyzed based on what it \emph{does} rather than what it \emph{says}. This approach is agnostic regarding phonetic values and focuses instead on the mathematical and logical structures that can be demonstrated through formal methods.

This is not a limitation but a methodological choice that enables rigorous, falsifiable claims within a well-defined scope.
